Om uit deze cyclus van bordelen en excessen te breken, is een verschuiving nodig in het bewustzijn van seksualiteit, menselijke waardigheid en verbinding. In plaats van het gebruik van bordelen als een remedie tegen emotionele eenzaamheid en onderdrukking, zouden gemeenschappen zich moeten richten op eerlijke en intieme relaties, educatie over seksuele gezondheid en een ruimte scheppen waarin gevoelens, verlangens en spanningen bespreekbaar zijn. Het elimineren van bordelen als noodzaak voor menselijke verbinding begint bij het veranderen van hoe men binnen de samenleving naar seksuele expressies kijkt, het wegnemen van stigma's en het creëren van veilige en respectvolle manieren voor mensen om hun innerlijke pijn en verlangens te helen. Sociale cohesie, persoonlijk contact en wederzijds respect kunnen de fundering zijn waarop een gezonder seksuele interactie rust.

Dus wat moet je doen als je worstelt met een pornoverslaving? Het eerste wat je moet doen, is jezelf echt deze vraag stellen: voegt porno iets toe aan mijn leven of neemt het iets weg? En je moet echt eerlijk zijn tegen jezelf over het antwoord op die vraag.
14:43
Dit is moeilijk voor sommigen van ons omdat als we graag naar porno kijken, we de gewoonte hebben om ons pornokijken te rechtvaardigen. We blijven onszelf vertellen dat iedereen het doet. Alleen jij kunt beslissen of pornografie je leven negatief beïnvloedt, maar totdat je besluit dat het je leven negatief beïnvloedt, zul je absoluut geen motivatie hebben om het te veranderen.
15:06
Stap twee: word je bewust. Je moet je volledig bewust worden van waar je probeert aan te ontsnappen door porno te gebruiken. De volgende keer dat je de drang voelt om ernaar te kijken, zit dan met de ongemakkelijke gevoelens alsof je de sensatie aan het verkennen bent en probeer te identificeren wat het is. Identificeer ook je triggers.
15:29
Een trigger is iets dat ervoor zorgt dat je de dwang voelt om naar porno te kijken. Denk aan de laatste uren of minuten die leiden tot het moment waarop de drang optreedt. Kun je triggers identificeren? Bijvoorbeeld, had je ruzie met iemand waardoor je stress voelde en dus ontlading nodig had?
15:45
Stap drie: ga in de richting van die negatieve gevoelstoestanden in plaats van ervan weg. De volgende keer dat je die drang voelt om porno te gebruiken, wil ik in plaats daarvan dat je afdaalt in dat gevoel, die negatieve gevoelstoestand, en ik wil dat je er onvoorwaardelijk aanwezig bij bent. Dit is het begin van integratie.
16:08
Om meer te begrijpen over hoe je dit proces kunt doen, wil ik dat je mijn video op YouTube bekijkt met de titel "How to heal the emotional body". Ook, als je merkt dat de gevoelstoestand die in feite je pornogebruik motiveert het gevoel van leegte is, wat zo vaak het geval is, wil ik dat je mijn YouTube-video bekijkt met de titel "How to stop feeling empty".
16:30
Vier: verander je leven in overeenstemming met de gevoelstoestand die je ontdekte dat je probeert te vermijden. Bijvoorbeeld, als je porno gebruikt om verveling te vermijden, word dan bezig. Vind een gezonde bezigheid die je gerichte aandacht en energie vereist. Of als je merkt dat je porno gebruikt om stress te ontvluchten, ontdek wat je het meest gestrest maakt en vind ofwel gezondere manieren om met de stress om te gaan of verander dat aspect van je leven volledig zodat de stress er niet meer is.
17:05
Kortom, we moeten het ding onder ogen zien en ermee werken waarvoor we porno gebruiken om te vermijden, om de oorzaak van ons pornogebruik te overwinnen. Pornogebruik is meer als een symptoom. Waarom zou je het symptoom aanpakken zonder de oorzaak aan te pakken?
17:19
Vijf: als je de drang voelt om naar porno te kijken, speel dan een spel van vooruitspoelen. Het is gemakkelijk om meegesleept te worden in de sensaties van opwinding en arousal, maar als je dat gevoel van opwinding kunt gebruiken om jezelf eraan te herinneren om de tape vooruit te spoelen naar na het kijken van porno - dat gevoel van teleurstelling, dat gevoel van uitputting, dat gevoel van lage eigenwaarde, het gevoel van totale en volledige mislukking en schaamte - dan kun je de herinnering aan hoe het zal zijn als je wel naar porno kijkt gebruiken om de intensiteit en de hitte van de opwinding te verminderen.
17:59
Zes: kom in de mentaliteit dat porno letterlijk niet bestaat. Als je de vorige stappen doet die ik net heb opgesomd, zal deze stap veel gemakkelijker zijn. Als het gaat om verslavingen, zoals ik eerder heb gezegd, is de verslaving op zich een symptoom van een dieper onderliggend probleem.
18:19
Dus wanneer we gewoonlijk verslavingen aanpakken, in plaats van ons te concentreren op de substantie waaraan we verslaafd zijn, willen we een manier vinden om de oorzaak van de verslaving op te lossen, om wat dan ook op te lossen dat de verslaving voedt, omdat als je dat doet, dan gaat de verslaving vanzelf weg.
18:40
Als het gaat om porno, kunnen we ons concentreren op het aanpakken van de oorzaak en ook het symptoom. We doen dit door cold turkey te gaan. Met andere woorden, als het gaat om porno, beperk je tijd niet tot het kijken naar porno, maar raak je uitlaatkleppen voor de verslaving kwijt, zodat je er letterlijk geen toegang toe hebt.
18:59
Denk en handel dan alsof porno niet bestaat. Uit pure noodzaak zul je gedwongen worden om de porno te vervangen door andere dingen om je energie op te richten.
19:10
Zeven: in lijn met het laatste punt dat ik maakte, wanneer je die drang voelt, herken die drang dan als een extreme uitbarsting van energie. Natuurlijk wil je die energie vrijlaten, maar de vraag is: naar wat? Wat ik wil dat je doet wanneer je de drang voelt om naar porno te kijken, is om die energie te herkanaliseren naar iets dat voordeliger is.
19:34
Je kunt je seksuele driften niet echt onderdrukken. Het is niet mogelijk. Het moet er op de een of andere manier uit. Het feit dat we onze seksuele driften niet kunnen onderdrukken, is in feite de reden waarom onze pornoabonnementen hoger zijn in staten zoals Utah, waar er een hoge mate van seksuele onderdrukking is. Onderdrukking werkt niet. Je moet die energie herkanaliseren naar iets anders.
19:55
Denk na over het benutten van die energie. Vraag jezelf af: als ik die energie zou kunnen benutten en inzetten voor iets anders dat mijn leven direct ten goede zou komen, wat zou dat dan zijn? En doe dat dan in plaats daarvan.
20:09
Als je dit genoeg doet, zul je beseffen dat je eigenlijk een leven hebt. Je kijkt niet langer naar porno, je voelt je verschrikkelijk over jezelf en je voelt dat je een leven moet krijgen.
20:21
Acht: verhoog je eigenwaarde en begin met zelfzorg. Zoals eerder besproken, vermindert porno onmiddellijk je eigenwaarde. De tegenovergestelde vibratie van porno is dus eigenwaarde. Als je verslaafd bent aan porno, is het een spiraal van zelfhaat geworden.
20:38
Om jezelf daaruit te trekken, begin dan gezond te leven en verander je perspectief zodat je jezelf in een positief licht kunt zien. Kijk naar de waarden die je wel hebt en de dingen waar je echt in gelooft, en begin volgens die waarden te leven. Leef vanuit een ruimte van integriteit met wie je werkelijk bent.
20:54
Doe dingen waar je van geniet. Eet een dieet dat je een goed gevoel over jezelf geeft. Beweeg. Verzorg jezelf op een manier die je trots maakt om jezelf aan de wereld te presenteren. Krijg hulp van buitenaf. Help anderen. Anderen helpen verhoogt je eigenwaarde vijfvoudig.
21:10
Zoek dingen die je aan het lachen maken. Begin met mediteren of ga naar yogales. Luister naar motiverende muziek die je inspireert om echt te leven. Begin je te concentreren op de dingen die je leuk vindt aan jezelf en vind goedkeuring voor de dingen die je niet leuk vindt aan jezelf. Dit zijn slechts enkele voorbeelden.
21:27
Verhoog je frequentie. Voor meer informatie over hoe je je frequentie kunt verhogen, kun je mijn YouTube-video opzoeken met de titel "How to raise your frequency and increase your vibration". Onthoud dat hoe meer je porno vervangt door voordeligere bezigheden, hoe beter en beter je leven zal worden.
21:40
Terwijl we het over dit onderwerp hebben, moet het bekend zijn dat als je denkt dat er iets niet goed is aan het kijken naar porno, je niet naar porno kunt kijken en je goed voelen.
21:53
Negen: nu ga ik je de belangrijkste tip geven als het gaat om elke verslaving, inclusief pornoverslaving. Lang geleden, toen wetenschappers studies deden met ratten en middelenverslaving, zouden ze hen in kooien stoppen en hun water mengen met middelen zoals cocaïne.
22:12
Uiteraard zouden deze ratten naar het water blijven gaan. Ze zouden verslaafd raken aan de cocaïne in het water en andere drugs tot op het punt dat ze zouden sterven. De voor de hand liggende aanname is dat de middelen de schuld waren voor de verslaving.
22:30
Toen merkte Bruce Alexander, die destijds een wetenschapper was, in 1970 iets interessants op. Deze ratten die verslaafd waren aan deze middelen, werden alleen in kooien geplaatst. Dus hij dacht bij zichzelf: wat is het volgende logische ding om te doen om te testen of middelen inderdaad de schuld zijn voor verslaving? Ah, laten we deze ratten in mooie kooien met vrienden plaatsen.
22:55
En wat ontdekte hij? Deze ratten raakten niet verslaafd aan de middelen. In feite leken ze het water dat was vermengd met deze specifieke medicijnen te mijden. Wat vertelt dit ons over verslaving? Het vertelt ons dat verslaving niet gaat over het middel; het gaat over een onderliggend gevoel van isolatie. De verslaving was een aanpassing aan een pijnlijke omgeving, vooral eentje die leeg was van verbinding.
23:26
Wat ik nu ga zeggen, wil ik dat je voor de rest van je leven onthoudt: emotionele isolatie is de nummer één oorzaak van alle verslavingen op deze planeet. Soberheid is niet het tegenovergestelde van verslaving; menselijke verbinding is dat. Dus maak verbinding met mensen.
23:44
Voor meer informatie over hoe je dit kunt doen, bekijk mijn YouTube-video met de titel "How to connect with someone".
23:48
Het is mijn wens dat seks niet wordt geassocieerd met zonde in de gedachten van mensen. We kunnen niet in oorlog zijn met onze seksualiteit en integrale en hele wezens worden in onszelf. Het is niet mijn wens dat porno de zondebok wordt voor de disfunctie van de samenleving, wat velen het zouden schilderen.
24:08
Ik ben er volledig voor om vrij en ruimdenkend te worden en om dingen te omarmen die nooit taboe zouden moeten zijn. Maar vrij en ruimdenkend zijn is niet hetzelfde als simpelweg het taboe vieren omdat het taboe is.
24:24
Ik sta open voor het idee dat wanneer mensen in een staat van afstemming komen met hun seksualiteit, waar hun relatie tot hun eigen seksualiteit gezond is en ook voordelig, dat we volledig in staat zijn om pornografie te maken die volkomen in balans is. Maar op dit moment is de overgrote meerderheid van porno diep uit balans. Het is een leeg pakket en je verdient iets substantieels.
